% ========================
% estilo.latex mínimo funcional
% ========================

\documentclass[12pt]{book} % report para capítulos

% ========================
% Paquetes y comandos extra
% ========================
% ===========================
% Paquetes básicos de idioma y codificación
% ===========================
\usepackage[utf8]{inputenc}   % Codificación UTF-8
\usepackage[T1]{fontenc}      % Acentos y caracteres correctos
\usepackage[spanish]{babel}   % Traducción al español (capítulos, índices, etc.)
\usepackage{csquotes}         % Citas tipográficas correctas

% ===========================
% Tipografía
% ===========================
\usepackage{lmodern}          % Fuente Latin Modern
\usepackage{microtype}        % Mejoras tipográficas (espaciado, justificación)

% ===========================
% Márgenes y geometría
% ===========================
\usepackage{geometry}         % Control de márgenes
\geometry{a4paper, top=3cm, bottom=3cm, left=3cm, right=3cm}

% ===========================
% Matemáticas
% ===========================
\usepackage{amsmath, amssymb, amsthm} % Paquetes AMS
\usepackage{mathtools}        % Extiende amsmath
\usepackage{physics}          % Notación física y matemática (derivadas, bra-ket, etc.)
\usepackage{siunitx}          % Unidades SI (e.g. \SI{3}{m/s})
% \sisetup{locale=ES}           % Configuración para español (coma decimal, etc.)
\AtBeginDocument{\RenewCommandCopy\qty\SI} % Resolve siunitx and physics conflict


% ===========================
% Gráficos, tablas y colores
% ===========================
\usepackage{graphicx}         % Insertar imágenes
\usepackage{xcolor}           % Colores personalizados
\usepackage{tikz}             % Dibujos vectoriales
\usetikzlibrary{calc,positioning,shapes,arrows} % Librerías útiles de TikZ
\usepackage{pgfplots}         % Gráficas de funciones
\pgfplotsset{compat=1.18}
\usepackage{float}            % Control de posición de figuras/tablas
\usepackage{booktabs}         % Tablas profesionales
\usepackage{multirow}         % Celdas que ocupan varias filas
\usepackage{array}            % Más control en tablas
\usepackage{colortbl}         % Tablas con colores
\usepackage{inconsolata}


% ===========================
% Listas y enumeraciones
% ===========================
\usepackage{enumitem}         % Control de listas enumeradas y viñetas

% ===========================
% Encabezados, pies y diseño
% ===========================
\usepackage{fancyhdr}         % Encabezados y pies de página
\usepackage{titlesec}         % Personalizar títulos de capítulos/secciones
\usepackage{setspace}         % Espaciado entre líneas
\usepackage{parskip}          % Control del espacio entre párrafos

% ===========================
% Referencias, hipervínculos y citas
% ===========================
\usepackage{hyperref}         % Hipervínculos en PDF
\hypersetup{
    colorlinks = true,
    linkcolor  = red!70,
    citecolor  = red!70,
    urlcolor   = red!70,
    pdfpagelayout = SinglePage, % Asegura que el contenido se ajuste a una sola página
    pdfstartview = Fit          % Ajusta el contenido al tamaño de la página
}
\usepackage{cleveref}         % Referencias inteligentes (\cref)

% ===========================
% Código fuente
% ===========================
\usepackage{listings}         % Mostrar código con estilo
\usepackage{minted}           % (mejor opción, requiere Python y pygments)

% ===========================
% Bibliografía
% ===========================
\usepackage[backend=biber,style=apa]{biblatex} % Ejemplo: estilo APA
\addbibresource{referencias.bib}              % Archivo .bib

% ===========================
% Otros útiles
% ===========================
\usepackage{pdfpages}         % Insertar PDFs externos
\usepackage{blindtext}        % Texto de prueba
\usepackage{caption}          % Personalizar pies de figura/tabla
\usepackage{subcaption}       % Subfiguras
\usepackage{tocloft} 
\usepackage{amsthm}
\usepackage{subcaption}
\usepackage{truncate} % permite truncar texto si no cabe
\usepackage{libertinus}  % reemplaza lmodern
\usepackage{booktabs}  % para \toprule, \midrule, \bottomrule
\usepackage{array}     % para definir columnas personalizadas
\usepackage{colortbl}  % colores en tablas
\usepackage{etoolbox}
\AtBeginEnvironment{tabular}{\rowcolors{2}{gray!10}{white}\renewcommand{\arraystretch}{1.2}}

% ===========================
% Opciones de fuentes sugeridas
% ===========================
% TeX Gyre Pagella (estilo Palatino)
% \usepackage{fontspec}
% \usepackage{unicode-math}
% \setmainfont{TeX Gyre Pagella}
% \setmathfont{TeX Gyre Pagella Math}

% TeX Gyre Termes (estilo Times)
% \setmainfont{TeX Gyre Termes}
% \setmathfont{TeX Gyre Termes Math}

% Libertinus (elegante y completa)
% \setmainfont{Libertinus Serif}
% \setmathfont{Libertinus Math}

% TeX Gyre Bonum (estilo Garamond)
% \setmainfont{TeX Gyre Bonum}
% \setmathfont{TeX Gyre Bonum Math}

% Latin Modern (moderno de Computer Modern)
% \setmainfont{Latin Modern Roman}
% \setmathfont{Latin Modern Math}


% \usepackage{helvet}
% \usepackage{libertine}
% \usepackage[sfdefault]{FiraSans}

\usepackage{tcolorbox} % para cajas de colores




  % si tienes paquetes personalizados
% aquí van los comandos personalizados
% Comando para incluir imágenes
\newcommand{\incluirimagen}[3][]{%
\begin{figure}[H]
    \centering
    \includegraphics[width=\linewidth,#1]{#2}
    \caption{#3}
    \label{fig:#2}
\end{figure}
}

% comando para ejercicios con fondo
\newtheoremstyle{ejerciciostyle}
  {10pt}   % Espacio arriba
  {10pt}   % Espacio abajo
  {\itshape} % Fuente del cuerpo
  {}       % Sangría
  {\bfseries} % Fuente del encabezado
  {}      % Puntuación tras encabezado
  { }      % Espacio tras encabezado
  {\thmname{#1} \thmnumber{#2}. \thmnote{#3}}


% % comando formal para enunciado de ejercicios
% \theoremstyle{ejerciciostyle}
% \newtheorem{ejercicio}{Ejercicio}[chapter]

\theoremstyle{ejerciciostyle}
\newtheorem{ejercicio}{Ejercicio}[section]

\renewcommand{\theejercicio}{\thechapter.\arabic{section}.\arabic{ejercicio}}


% comando formal para soluciones
\theoremstyle{remark}
\newtheorem*{solucion}{Solución}

% Comando para dos imágenes en paralelo
\newcommand{\dosimagenes}[6]{%
    \begin{figure}[h!]
        \centering
        \begin{minipage}{0.48\linewidth}
            \centering
            \includegraphics[width=\linewidth]{#1}
            \caption{#2}
            \label{#5}
        \end{minipage}\hfill
        \begin{minipage}{0.48\linewidth}
            \centering
            \includegraphics[width=\linewidth]{#3}
            \caption{#4}
            \label{#6}
        \end{minipage}
    \end{figure}
}

% \dosimagenes{media/fondo.jpg}{Descripción 1}{media/fondo.jpg}{Descripción 2}{fig:descripcion1}{fig:descripcion2}

% \ref{fig:descripcion1} es la mejor
% \ref{fig:descripcion2} es la mejor

\newcommand{\portadaimg}{\VAR{portadaimg}}

% Comando para crear una nota estilo información
\newcommand{\nota}[2]{%
\begin{tcolorbox}[colframe=blue!75!black, colback=blue!5!white, title=\textbf{#1}]
    #2
\end{tcolorbox}
}
  % comandos LaTeX propios
% ===========================
% Diseño general
% ===========================
\setstretch{1.15} % interlineado
\setlength{\parskip}{0.5em} % espacio entre párrafos
\setlength{\parindent}{0pt} % sin sangría

% ===========================
% Estilo de capítulos y secciones (titlesec)
% ===========================
\titleformat{\chapter}[display]
  {\bfseries\Huge}
  {\filleft\Large\scshape Capítulo \thechapter}
  {1ex}
  {\titlerule[1pt]\vspace{1ex}\filright}
  [\vspace{1ex}\titlerule]

\titlespacing*{\chapter}{0pt}{0pt}{2em}

\titleformat{\section}
  {\Large\bfseries}
  {\thesection}{1em}{}

\titleformat{\subsection}
  {\large\bfseries}
  {\thesubsection}{1em}{}

\titleformat{\subsubsection}
  {\normalsize\bfseries\itshape}
  {\thesubsubsection}{1em}{}

% ===========================
% Encabezados y pies de página (fancyhdr)
% ===========================
\pagestyle{fancy}
\fancyhf{} % limpia
\fancyhead[L]{\small\scshape\nouppercase{\leftmark}} % sección/capítulo en mayúsculas pequeñas
\fancyhead[R]{\small\thepage}                        % número de página
%\fancyfoot[C]{\scriptsize\itshape Apuntes de la carrera} % texto fijo abajo en cursiva
% Encabezados y pies de página personalizados
% \fancyfoot[L]{\scriptsize\itshape Nombre de la asignatura} % pie de página izquierdo en cursiva
\fancyfoot[R]{\scriptsize\itshape Ismael Sallami Moreno}        % pie de página derecho con el nombre del autor

% Línea bajo el encabezado
\renewcommand{\headrulewidth}{0.5pt} % línea más gruesa en el encabezado
% Línea en el pie
\renewcommand{\footrulewidth}{0.4pt} % línea fina en el pie
\renewcommand{\sectionmark}[1]{%
  \markboth{\thesection\quad #1}{}%
}

% ===========================
% Numeración de elementos
% ===========================
\numberwithin{equation}{chapter} % ecuaciones numeradas por capítulo
\numberwithin{figure}{chapter}   % figuras numeradas por capítulo
\numberwithin{table}{chapter}    % tablas numeradas por capítulo

% ===========================
% Listas y enumeraciones
% ===========================
\setlist[itemize]{label=--, left=1.5em}
\setlist[enumerate]{label=\arabic*), left=1.5em}

% ===========================
% Estilo de citas y bibliografía
% ===========================
\DefineBibliographyStrings{spanish}{%
  references = {Bibliografía},
}

% ===========================
% Entornos personalizados
% ===========================
\newtheoremstyle{cajita} % nombre del estilo
  {1em}   % espacio arriba
  {1em}   % espacio abajo
  {}      % fuente del cuerpo
  {}      % indentación
  {\bfseries} % fuente del título
  {.}     % puntuación tras título
  {0.5em} % espacio tras título
  {\thmname{#1}\thmnumber{ #2} \thmnote{(#3)}} % formato

\theoremstyle{cajita}
\newtheorem{teorema}{Teorema}[chapter]
\newtheorem{definicion}{Definición}[chapter]
\newtheorem{ejemplo}{Ejemplo}[chapter]
\newtheorem{proposicion}{Proposición}[chapter]

% ===========================
% Configuración de lstlisting
% ===========================

% ===============================================
% ESTILO 1: MODERNO Y MINIMALISTA
% ===============================================

% Definir colores personalizados
\definecolor{codegreen}{rgb}{0,0.6,0}
\definecolor{codegray}{rgb}{0.5,0.5,0.5}
\definecolor{codepurple}{rgb}{0.58,0,0.82}
\definecolor{backcolour}{rgb}{0.95,0.95,0.92}
\definecolor{framecolor}{rgb}{0.8,0.8,0.8}

\lstset{
    backgroundcolor=\color{backcolour},   
    commentstyle=\color{codegreen},
    keywordstyle=\color{magenta},
    numberstyle=\tiny\color{codegray},
    stringstyle=\color{codepurple},
    basicstyle=\ttfamily\footnotesize,
    breakatwhitespace=false,         
    breaklines=true,                 
    captionpos=b,                    
    keepspaces=true,                 
    numbers=left,                    
    numbersep=5pt,                  
    showspaces=false,                
    showstringspaces=false,
    showtabs=false,                  
    tabsize=2,
    frame=shadowbox,
    frameround=tttt,
    rulecolor=\color{framecolor},
    rulesepcolor=\color{framecolor},
    xleftmargin=20pt,
    xrightmargin=20pt,
    aboveskip=20pt,
    belowskip=20pt
}

% ===============================================
% ESTILO 2: ELEGANTE CON BORDES REDONDEADOS
% ===============================================

% Colores para estilo elegante
\definecolor{lightblue}{rgb}{0.93,0.95,1}
\definecolor{darkblue}{rgb}{0.1,0.2,0.5}
\definecolor{mediumblue}{rgb}{0.2,0.4,0.8}
\definecolor{darkgreen}{rgb}{0,0.5,0}
\definecolor{darkred}{rgb}{0.6,0,0}

\lstdefinestyle{elegant}{
    backgroundcolor=\color{lightblue},
    commentstyle=\color{darkgreen}\itshape,
    keywordstyle=\color{darkblue}\bfseries,
    numberstyle=\tiny\color{gray},
    stringstyle=\color{darkred},
    basicstyle=\ttfamily\small,
    breakatwhitespace=false,
    breaklines=true,
    captionpos=t,
    keepspaces=true,
    numbers=left,
    numbersep=8pt,
    showspaces=false,
    showstringspaces=false,
    showtabs=false,
    tabsize=4,
    frame=single,
    frameround=tttt,
    framesep=10pt,
    xleftmargin=15pt,
    xrightmargin=15pt,
    aboveskip=15pt,
    belowskip=15pt,
    columns=flexible
}

% ===============================================
% ESTILO 3: PROFESIONAL CORPORATIVO
% ===============================================

% Colores corporativos
\definecolor{corporatebg}{rgb}{0.98,0.98,0.98}
\definecolor{corporateblue}{rgb}{0.07,0.29,0.49}
\definecolor{corporategray}{rgb}{0.4,0.4,0.4}
\definecolor{corporategreen}{rgb}{0.13,0.55,0.13}
\definecolor{corporatered}{rgb}{0.8,0.2,0.2}

\lstdefinestyle{corporate}{
    backgroundcolor=\color{corporatebg},
    commentstyle=\color{corporategreen}\slshape,
    keywordstyle=\color{corporateblue}\bfseries,
    numberstyle=\scriptsize\color{corporategray},
    stringstyle=\color{corporatered},
    basicstyle=\ttfamily\footnotesize,
    breakatwhitespace=false,
    breaklines=true,
    captionpos=b,
    keepspaces=true,
    numbers=left,
    numbersep=12pt,
    showspaces=false,
    showstringspaces=false,
    showtabs=false,
    tabsize=3,
    frame=leftline,
    framerule=3pt,
    rulecolor=\color{corporateblue},
    xleftmargin=25pt,
    aboveskip=20pt,
    belowskip=20pt,
    lineskip=1pt
}

% ===============================================
% ESTILO 4: MODERNO CON SOMBRAS
% ===============================================

% Colores modernos
\definecolor{modernbg}{rgb}{0.97,0.97,0.97}
\definecolor{moderngray}{rgb}{0.3,0.3,0.3}
\definecolor{modernpurple}{rgb}{0.5,0.2,0.8}
\definecolor{modernteal}{rgb}{0,0.5,0.5}
\definecolor{modernorange}{rgb}{0.8,0.4,0}

\lstdefinestyle{modern}{
    backgroundcolor=\color{modernbg},
    commentstyle=\color{modernteal}\itshape,
    keywordstyle=\color{modernpurple}\bfseries,
    numberstyle=\tiny\color{moderngray},
    stringstyle=\color{modernorange},
    basicstyle=\ttfamily\small,
    breakatwhitespace=false,
    breaklines=true,
    captionpos=t,
    keepspaces=true,
    numbers=left,
    numbersep=10pt,
    showspaces=false,
    showstringspaces=false,
    showtabs=false,
    tabsize=4,
    frame=tb,
    framerule=2pt,
    rulecolor=\color{modernpurple},
    xleftmargin=20pt,
    xrightmargin=20pt,
    aboveskip=25pt,
    belowskip=25pt
}

% ===============================================
% CONFIGURACIÓN PARA DIFERENTES LENGUAJES
% ===============================================

% Python
\lstdefinestyle{python}{
    language=Python,
    style=elegant,
    morekeywords={True,False,None,self,cls,def,class,import,from,as,with,yield,async,await},
    morecomment=[l]{\#},
    morestring=[b]',
    morestring=[b]"
}

% Java
\lstdefinestyle{java}{
    language=Java,
    style=corporate,
    morekeywords={var,record,sealed,permits,non-sealed}
}

% C++
\lstdefinestyle{cpp}{
    language=C++,
    style=modern,
    morekeywords={constexpr,nullptr,auto,decltype,override,final}
}

% JavaScript
\lstdefinestyle{javascript}{
    language=Java,
    style=elegant,
    morekeywords={let,const,var,function,class,extends,import,export,default,async,await,yield},
    morecomment=[l]{//},
    morecomment=[s]{/*}{*/},
    morestring=[b]',
    morestring=[b]",
    morestring=[b]`
}

% ===============================================
% EJEMPLOS DE USO
% ===============================================

% Para usar el estilo por defecto:
% \begin{lstlisting}
% código aquí
% \end{lstlisting}

% Para usar un estilo específico:
% \begin{lstlisting}[style=elegant]
% código aquí
% \end{lstlisting}

% Para incluir un archivo con estilo específico:
% \lstinputlisting[style=python]{archivo.py}

% Para código inline:
% \lstinline[style=modern]{código inline}

% ===============================================
% CONFIGURACIÓN ADICIONAL PARA TÍTULOS Y CARACTERES
% ===============================================

% Personalizar el formato de los títulos de los listados
\renewcommand\lstlistingname{Código}
\renewcommand\lstlistlistingname{Lista de Códigos}

% Configurar el formato del título con soporte para tildes
\lstset{
    title=\lstname,
    captionpos=t,
    abovecaptionskip=10pt,
    belowcaptionskip=5pt,
    % Configuración global para caracteres especiales
    inputencoding=utf8,
    extendedchars=true
}

% ===============================================
% COMANDOS PERSONALIZADOS ÚTILES
% ===============================================

% Comando para código inline con soporte automático de tildes
\newcommand{\codeinline}[2][modern]{\lstinline[style=#1,inputencoding=utf8,extendedchars=true]{#2}}

% Comando para bloques de código con título personalizado
\newcommand{\codeblock}[3][elegant]{%
    \begin{lstlisting}[style=#1,caption={#2},label={lst:#2},inputencoding=utf8,extendedchars=true]
    #3
    \end{lstlisting}
}

% Comando para incluir archivos con configuración automática
\newcommand{\includecode}[3][python]{%
    \lstinputlisting[style=#1,caption={#3},label={lst:#3},inputencoding=utf8,extendedchars=true]{#2}
}

% ===============================================
% CONFIGURACIONES ESPECIALES PARA IDIOMAS
% ===============================================

% Configuración específica para código en español
\lstdefinestyle{español}{
    style=elegant,
    inputencoding=utf8,
    extendedchars=true,
    % Palabras clave en español para pseudocódigo
    morekeywords={función,procedimiento,inicio,fin,si,entonces,sino,mientras,para,hasta,hacer,repetir,caso,segun,verdadero,falso,entero,real,caracter,cadena,booleano,leer,escribir,imprimir}
}

% Configuración para comentarios multilíngües
\lstset{
    morecomment=[l]{//\ },
    morecomment=[l]{\#\ },
    morecomment=[s]{/*}{*/},
    morecomment=[s]{<!--}{-->}
}

% ===============================================
% CONFIGURACIÓN PARA DIFERENTES LENGUAJES
% ===============================================

% Python
\lstdefinestyle{style1}{
    language=Python,
    style=elegant,
    morekeywords={True,False,None,self,cls,def,class,import,from,as,with,yield,async,await},
    morecomment=[l]{\#},
    morestring=[b]',
    morestring=[b]",
    % Soporte para caracteres especiales
    inputencoding=utf8,
    extendedchars=true
}

% Java
\lstdefinestyle{style2}{
    language=Java,
    style=corporate,
    morekeywords={var,record,sealed,permits,non-sealed},
    % Soporte para caracteres especiales
    inputencoding=utf8,
    extendedchars=true
}

% C++
\lstdefinestyle{style3}{
    language=C++,
    style=modern,
    morekeywords={constexpr,nullptr,auto,decltype,override,final},
    % Soporte para caracteres especiales
    inputencoding=utf8,
    extendedchars=true
}

% ===========================
% Estilo global de tablas
% ===========================

\usepackage{booktabs}   % reglas profesionales
\usepackage{colortbl}   % color en filas
\usepackage{xcolor}     % colores
\usepackage{float}      % [H]

% Color de filas alternadas
\rowcolors{2}{gray!10}{white}

% Espacio vertical entre filas
\renewcommand{\arraystretch}{1.2}

% Cambiar el tamaño de columna por defecto
\setlength{\tabcolsep}{8pt}

% Redefinir tabla para que todas las tablas tengan el estilo
\let\oldtabular\tabular
\let\endoldtabular\endtabular
\renewenvironment{tabular}[1]{%
  \oldtabular{#1}%
}{%
  \endoldtabular
}
\usepackage{longtable,booktabs,xcolor}
\rowcolors{2}{gray!10}{white}   % filas alternadas
\renewcommand{\arraystretch}{1.2} % espacio vertical entre filas







   % estilos de secciones, etc.

% ========================
% Configuración índice y listas
% ========================
\setlength{\cftbeforesecskip}{5pt}
\setlength{\headheight}{14pt}  % un poco más que 13.6pt

\renewcommand{\normalsize}{\fontsize{10}{12}\selectfont}

% Fix para listas de Pandoc
\providecommand{\tightlist}{%
  \setlength{\itemsep}{0pt}\setlength{\parskip}{0pt}}



%===============
% ESPACIOS
%===============

% --- Compactar secciones ---
\titlespacing*{\section}{0pt}{1.2ex plus 0.5ex minus 0.2ex}{0.8ex}
\titlespacing*{\subsection}{0pt}{1ex plus 0.3ex minus 0.2ex}{0.5ex}

% --- Compactar flotantes (figuras/tablas) ---
\setlength{\textfloatsep}{8pt}
\setlength{\intextsep}{6pt}
\setlength{\floatsep}{6pt}

% --- Compactar listas ---
\setlist{nosep}

% --- Espacio entre párrafos ---
\setlength{\parskip}{4pt}



%=======================
% fancy with parameters
%=======================
%\fancyfoot[L]{\scriptsize\itshape Modelos de Computación}
\fancyfoot[L]{\normalsize Modelos de
Computación} % pie de página izquierdo con tamaño normal

\setcounter{tocdepth}{1} % Muestra solo hasta subsecciones en el índice

% ========================
% Inicio del documento
% ========================
\begin{document}

% Cambiar puntos suspensivos en el índice
\renewcommand{\cftsecleader}{\cftdotfill{\cftdotsep}}

% Ajustar formato de secciones y subsecciones en el índice
\renewcommand{\cftsecfont}{\bfseries} % Secciones en negrita
\renewcommand{\cftsecpagefont}{\bfseries} % Números de página en negrita para secciones
\renewcommand{\cftsubsecfont}{\normalfont} % Subsecciones en formato normal
\renewcommand{\cftsubsecpagefont}{\normalfont} % Números de página en formato normal para subsecciones

% Espaciado entre entradas del índice
\setlength{\cftbeforesecskip}{8pt} % Espaciado antes de secciones
\setlength{\cftbeforesubsecskip}{4pt} % Espaciado antes de subsecciones



%% portada.tex sencilla
\begin{titlepage}
\newgeometry{top=2cm,bottom=2cm,left=2.5cm,right=2.5cm}

\begin{center}

% Imagen de portada (portadaimg viene de Pandoc -M portadaimg="...")
\includegraphics[width=\textwidth]{$portadaimg$}

\vspace{2cm}

% Nombre de la asignatura (también desde Pandoc)
{\Huge \bfseries $asignatura$ \par}

\vfill

{\large Autor: \textbf{Ismael Sallami Moreno} \par}
\vspace{0.5cm}
{\large \today}

\end{center}

\restoregeometry
\end{titlepage}



%==========================
% PORTADA: ENTRADA MANUAL
%==========================

% portada.tex
\begin{titlepage}
    \newgeometry{top=2cm,bottom=2cm,left=2.5cm,right=2.5cm} % márgenes personalizados
    
    % Fondo con transparencia
    \begin{tikzpicture}[remember picture,overlay]
        % \node[opacity=0.15,inner sep=0pt] at (current page.center)
        \node[inner sep=0pt] at (current page.center)
            {\includegraphics[width=\paperwidth,height=\paperheight]{../../../extraFiles/img/fondo_info.jpg}};
    \end{tikzpicture}

    % Contenido de la portada
    \begin{center}
        \vspace*{2cm}
        
        \vspace{5cm} % Añadir más espacio antes del contenido
        {\Huge \bfseries\scshape\textcolor{white}{Modelos de
Computación} \par}
        \vspace{0.5cm}
        {\Large \itshape\textcolor{white}{Temario} \par}
        \vspace{0.5cm}
        % {\small \itshape \href{https://ismael-sallami.github.io}{https://ismael-sallami.github.io} \par}
        % {\small \itshape \href{https://elblogdeismael.github.io}{https://elblogdeismael.github.io} \par}


        \vfill
        
        % {\LARGE Ismael Sallami Moreno \par}

        \begin{flushright}
            {Ismael Sallami Moreno \par}
            {\small \itshape \href{https://elblogdeismael.github.io}{Recursos Ingeniería Informática y Ade} \par}
        \end{flushright}
        \vspace{0.3cm}
        % {\Large Universidad de Granada \par}
        
        % \vspace{1cm}
        % \includegraphics[width=0.25\textwidth]{../../../extraFiles/img/ugr.png} % opcional: logo
        % \vspace{1cm}
        
        % {\large \today}
    \end{center}
    
    \restoregeometry
\end{titlepage}


%==========================
% LICENCIA
%==========================

\begin{tikzpicture}[remember picture,overlay]
\node[anchor=south west, xshift=1cm, yshift=1cm] at (current page.south west) {
\begin{minipage}{0.4\textwidth}
\begin{flushleft}
\section*{Licencia}

Este trabajo está bajo una 
\href{https://creativecommons.org/licenses/by-nc-nd/4.0/}{Licencia Creative Commons BY-NC-ND 4.0}.

\bigskip

Permisos: Se permite compartir, copiar y redistribuir el material en cualquier medio o formato.

\bigskip

Condiciones: Es necesario dar crédito adecuado, proporcionar un enlace a la licencia e indicar si se han realizado cambios. No se permite usar el material con fines comerciales ni distribuir material modificado.

\bigskip

\begin{center}
  \href{https://creativecommons.org/licenses/by-nc-nd/4.0/}{\includegraphics[width=0.35\textwidth]{../../../extraFiles/img/by-nc-nd.png}}
\end{center}
\end{flushleft}
\end{minipage}
};
\end{tikzpicture}

\thispagestyle{empty}
\clearpage

%==========================
% AUTOR
%==========================

% Página del autor
\begin{center}
    \vspace*{3cm} % Añadir más espacio en la parte superior
    {\Huge Modelos de
Computación}\\[2cm] % Incrementar el espacio entre líneas
    {\Large Ismael Sallami Moreno}\\[1cm] % Incrementar el espacio entre líneas
    % \includegraphics[width=0.3\textwidth]{autor.jpg} % opcional foto
    \vfill % Añadir espacio flexible para centrar verticalmente
\end{center}
\thispagestyle{empty}

% Página en blanco sin estilo
\newpage
\thispagestyle{empty}
\mbox{}

% Otra página en blanco sin estilo
\newpage
\thispagestyle{empty}
\mbox{}


%==========================
% BIOGRAFÍA
%==========================

% Breve descripción del autor
% \chapter*{Biografía}
% \addcontentsline{toc}{chapter}{Biografía} % aparece en índice
% Aquí escribes una breve descripción sobre ti, tu formación, experiencia, etc.
% \cleardoublepage


% % ===============================
% licencia.tex
% ===============================
\section*{Licencia}

Este trabajo está licenciado bajo una 
\href{https://creativecommons.org/licenses/by-nc-nd/4.0/}{Licencia Creative Commons Reconocimiento-NoComercial-SinObraDerivada 4.0 Internacional}.

\bigskip

\textbf{Usted es libre de:}
\begin{itemize}
  \item \textbf{Compartir} — copiar y redistribuir el material en cualquier medio o formato.
\end{itemize}

\bigskip

\textbf{Bajo los siguientes términos:}
\begin{description}
  \item[\textbf{Reconocimiento}] Debe otorgar el crédito adecuado, proporcionar un enlace a la licencia e indicar si se han realizado cambios. Puede hacerlo de cualquier manera razonable, pero no de una manera que sugiera que tiene el apoyo del licenciante o lo recibe por el uso que hace.

  \item[\textbf{NoComercial}] No puede utilizar el material para fines comerciales.

  \item[\textbf{SinObraDerivada}] Si remezcla, transforma o crea a partir del material, no puede distribuir el material modificado.
\end{description}

\bigskip

\begin{center}
  \href{https://creativecommons.org/licenses/by-nc-nd/4.0/}{\includegraphics[width=0.35\textwidth]{../../../extraFiles/img/by-nc-nd.png}}
\end{center}
  % licencia
% \thispagestyle{empty} % quitar número de página en la portada
% \clearpage

% --- Índice ---
\tableofcontents
% \listoffigures
\clearpage

%\listoftables
%\clearpage
%\thispagestyle{empty} % quitar número de página en la portada
%\clearpage
%
% Índice de código
%\renewcommand{\lstlistlistingname}{Índice de Código}
%\lstlistoflistings
%\clearpage
%
% Índice de ecuaciones
%\renewcommand{\listtheoremname}{Índice de Ecuaciones}
%\listoftheorems[ignoreall,show={equation}]
%\clearpage

% --- Contenido Markdown generado por Pandoc ---
\part{Teoría}
\chapter{Introducción a la Computación}

\section{Fundamentos y Orígenes}

\subsection{Preguntas Clave: Definición, problemas resolubles, y Modelos de Computación}

La Ciencia de la Computación se fundamenta en la búsqueda de respuestas a preguntas trascendentales sobre los límites de la resolución automática de problemas. Entre las cuestiones centrales se encuentran:
\begin{itemize}
    \item \textbf{¿Qué puede ser resuelto de forma automática?} Esta pregunta indaga sobre la existencia de algoritmos para problemas dados, es decir, procedimientos mecánicos sistemáticos que garantizan una terminación para cualquier entrada válida.
    \item \textbf{¿Qué puede ser resuelto de forma eficiente?} Más allá de la mera resolubilidad, este interrogante aborda la viabilidad práctica de los algoritmos en términos de recursos computacionales como tiempo y espacio.
    \item \textbf{¿Qué estructuras son comunes en la computación con símbolos y cómo pueden ser procesadas?} Esta línea de investigación se enfoca en el estudio de patrones en conjuntos de cadenas y en los mecanismos formales para su manipulación.
\end{itemize}
Para abordar estas preguntas, se desarrollan \textbf{Modelos de Computación}. Estos modelos son abstracciones matemáticas de un dispositivo de cómputo, que permiten analizar sus capacidades y limitaciones de manera formal. Un concepto clave en este ámbito es el de \textbf{autómata}, concebido como un sistema algebraico que procesa información. Un autómata, en su forma más general, se define por sus conjuntos de estados, señales de entrada y señales de salida, junto con las operaciones que gobiernan sus transiciones y respuestas. El estudio de estos modelos, como las Máquinas de Turing y sus simplificaciones (autómatas finitos, autómatas con pila), constituye la base de la Teoría de la Computabilidad y de Lenguajes Formales.

\subsection{Breve Historia: Desde precursores hasta Autómatas y el Problema de la Parada}

Los fundamentos de la computación se remontan a los trabajos de lógicos y matemáticos como Russell, Hilbert y Boole, quienes sentaron las bases del formalismo matemático moderno. En las décadas de 1930 y 1940, figuras como Alan Turing y Alonzo Church formalizaron la noción de "computabilidad". Turing, en particular, propuso un modelo teórico, la \textbf{Máquina de Turing}, que se considera la definición de un dispositivo de cómputo universal.

En su búsqueda por resolver los 23 problemas propuestos por David Hilbert, la comunidad científica demostró que ciertos problemas son \textbf{indecidibles}, es decir, no admiten una solución algorítmica general. El ejemplo canónico es el \textbf{Problema de la Parada (Halting Problem)}.

\begin{definicion}[Problema de la Parada]
    No existe un programa de ordenador universal, llamémoslo \texttt{Stops(P, x)}, que pueda tomar como entrada cualquier otro programa \texttt{P} y unos datos \texttt{x} y determinar, en un tiempo finito, si \texttt{P} terminará su ejecución con la entrada \texttt{x} o si entrará en un bucle infinito.
\end{definicion}

La imposibilidad de resolver este problema fue un resultado trascendental que estableció los límites fundamentales de lo que puede ser computado.

Paralelamente, el desarrollo de los primeros lenguajes de programación como FORTRAN, COBOL y LISP en los años 50 impulsó el estudio de la sintaxis. En esta década, Noam Chomsky, junto con Michael Rabin y Dana Scott, formalizó la teoría de los \textbf{Autómatas y Lenguajes Formales}, estableciendo una jerarquía de gramáticas y máquinas que reconocen distintas clases de patrones.

\section{Elementos Básicos: Lenguajes Formales}

\subsection{Alfabetos}

La teoría de lenguajes formales se construye a partir de conceptos básicos. El más fundamental es el de alfabeto.

\begin{definicion}[Alfabeto]
Un \textbf{alfabeto} es un conjunto \textbf{finito} y no vacío de elementos denominados \textbf{símbolos} o \textbf{letras}.
\end{definicion}

\begin{anotacion}
Se denotan los alfabetos con letras mayúsculas como $\Sigma$, $A$, $B$, etc., y sus símbolos con letras minúsculas como $a, b, c$ o números.
\end{anotacion}

\begin{ejemplo}
Algunos alfabetos comunes son:
\begin{itemize}
    \item El alfabeto binario: $\Sigma = \{0, 1\}$.
    \item El alfabeto de letras minúsculas: $\Sigma = \{a, b, c, \dots, z\}$.
    \item Un alfabeto de vectores: $B = \{\langle 0,0 \rangle, \langle 0,1 \rangle, \langle 1,0 \rangle, \langle 1,1 \rangle\}$.
\end{itemize}
\end{ejemplo}

\subsection{Palabras}

\begin{definicion}[Palabra]
Una \textbf{palabra} (también llamada \textbf{cadena} o \textbf{string}) sobre un alfabeto $A$ es una sucesión finita de símbolos de $A$. La \textbf{longitud} de una palabra $u$, denotada por $|u|$, es el número de símbolos que la componen.
\end{definicion}

\begin{anotacion}
Las palabras se denotan con letras minúsculas como $u, v, w, x, y, z$. El conjunto de todas las palabras posibles sobre un alfabeto $A$ se denota como $A^*$.
\end{anotacion}

\begin{definicion}[Palabra Vacía]
La \textbf{palabra vacía}, denotada por $\epsilon$ (o a veces $\lambda$), es la única palabra de longitud cero, $| \epsilon | = 0$. Es un elemento de $A^*$ para cualquier alfabeto $A$.
\end{definicion}

\begin{anotacion}
El conjunto de todas las palabras no vacías sobre un alfabeto $A$ se denota como $A^+$. Por lo tanto, $A^+ = A^* \setminus \{\epsilon\}$.
\end{anotacion}

El conjunto de palabras $A^*$ con la operación de concatenación forma una estructura algebraica conocida como monoide, donde la palabra vacía $\epsilon$ es el elemento neutro.

\subsection{Lenguajes}

\begin{definicion}[Lenguaje]
Un \textbf{lenguaje} $L$ sobre un alfabeto $A$ es cualquier subconjunto de $A^*$, es decir, $L \subseteq A^*$. Las palabras que pertenecen al lenguaje se denominan a menudo "cadenas del lenguaje" o "frases".
\end{definicion}

\begin{ejemplo}
    Dado el alfabeto $A = \{a, b\}$:
    \begin{itemize}
        \item $L_1 = \{a, b, \epsilon\}$ es un lenguaje finito con tres palabras.
        \item $L_2 = \{a^ib^i \mid i \geq 0\}$ es un lenguaje infinito que contiene palabras con igual número de a's seguidas de b's.
        \item $L_3 = \{uu^{-1} \mid u \in \{a, b\}^*\}$ es el lenguaje de los palíndromos de longitud par.
        \item $L_4 = \{a^{n^2} \mid n \geq 1\}$ es el lenguaje de cadenas de a's cuya longitud es un cuadrado perfecto.
    \end{itemize}
\end{ejemplo}

Una propiedad fundamental del conjunto de todos los lenguajes sobre un alfabeto no vacío es su \textbf{no numerabilidad}. Mientras que el conjunto de todas las palabras $A^*$ es numerable, el conjunto de todos sus subconjuntos (es decir, el conjunto de todos los lenguajes, $\mathcal{P}(A^*)$) es no numerable. Esto tiene una profunda implicación: existen más lenguajes que programas para describirlos o reconocerlos, lo que demuestra que debe haber problemas (lenguajes) que no son computacionalmente resolubles.

\section{Operaciones Esenciales}

\subsection{Sobre Palabras}

Existen varias operaciones fundamentales que se pueden aplicar a las palabras.

\begin{definicion}[Concatenación]
Dados $u = a_1 \dots a_n$ y $v = b_1 \dots b_m$ dos palabras sobre un alfabeto $A$, su \textbf{concatenación}, denotada $uv$, es la palabra $a_1 \dots a_n b_1 \dots b_m$.
\end{definicion}
La concatenación es asociativa y tiene a $\epsilon$ como elemento neutro ($u\epsilon = \epsilon u = u$).

\begin{definicion}[Iteración]
La \textbf{iteración n-ésima} de una palabra $u$, denotada $u^n$, es la concatenación de $u$ consigo misma $n$ veces. Se define formalmente como $u^0 = \epsilon$ y $u^{i+1} = u^i u$ para $i \geq 0$.
\end{definicion}

\begin{definicion}[Palabra Inversa]
La \textbf{palabra inversa} de $u = a_1 \dots a_n$, denotada $u^{-1}$, es la palabra $a_n \dots a_1$.
\end{definicion}

\begin{anotacion}
La operación de inversión no es un homomorfismo. Por ejemplo, si $f(u) = u^{-1}$, entonces $f(uv) = (uv)^{-1} = v^{-1}u^{-1}$, que en general es distinto de $f(u)f(v) = u^{-1}v^{-1}$.
\end{anotacion}

\subsection{Sobre Lenguajes}

Al ser los lenguajes conjuntos de palabras, heredan las operaciones de conjuntos como la \textbf{unión} ($L_1 \cup L_2$), \textbf{intersección} ($L_1 \cap L_2$) y \textbf{complementario} ($\bar{L} = A^* \setminus L$). Adicionalmente, se definen operaciones específicas.

\begin{definicion}[Concatenación de Lenguajes]
La \textbf{concatenación} de dos lenguajes $L_1$ y $L_2$ es el lenguaje $L_1 L_2 = \{uv \mid u \in L_1, v \in L_2\}$.
\end{definicion}
Esta operación es asociativa y tiene como elemento neutro el lenguaje $\{\epsilon\}$.

\begin{definicion}[Clausura de Kleene y Clausura Positiva]
Dado un lenguaje $L$, su \textbf{clausura de Kleene} (o estrella), denotada $L^*$, es la unión de todas sus potencias:
\[ L^* = \bigcup_{i \geq 0} L^i = L^0 \cup L^1 \cup L^2 \cup \dots \]
La \textbf{clausura positiva}, denotada $L^+$, se define de manera similar pero excluyendo la potencia cero:
\[ L^+ = \bigcup_{i \geq 1} L^i = L^1 \cup L^2 \cup \dots \]
donde $L^0 = \{\epsilon\}$ y $L^{i+1} = L L^i$.
\end{definicion}
Se cumple que $L^+ = L^*$ si y solo si $\epsilon \in L$. Si $\epsilon \notin L$, entonces $L^* = L^+ \cup \{\epsilon\}$. Si $L$ es un lenguaje no vacío, $L^*$ siempre es infinito.

\section{Generación de Lenguajes}

\subsection{Gramáticas: Definición y Lenguaje Generado}
Los lenguajes, especialmente los infinitos, requieren un mecanismo finito para su descripción. Las gramáticas generativas, introducidas por Noam Chomsky, son uno de dichos mecanismos.

\begin{definicion}[Gramática Generativa]
Una \textbf{gramática generativa} es una tupla $G = (V, T, P, S)$, donde:
\begin{itemize}
    \item $V$ es un alfabeto finito de \textbf{variables} o símbolos no terminales.
    \item $T$ es un alfabeto finito de \textbf{símbolos terminales}, disjunto de $V$.
    \item $P$ es un conjunto finito de \textbf{reglas de producción} de la forma $\alpha \rightarrow \beta$, donde $\alpha, \beta \in (V \cup T)^*$ y $\alpha$ contiene al menos una variable.
    \item $S \in V$ es el \textbf{símbolo inicial} o axioma.
\end{itemize}
\end{definicion}

Una palabra $\beta$ se \textbf{deriva} de $\alpha$ en un paso ($\alpha \Rightarrow \beta$) si existe una regla $\gamma \rightarrow \phi \in P$ y $\alpha$ puede escribirse como $\alpha_1\gamma\alpha_2$, de tal forma que $\beta = \alpha_1\phi\alpha_2$. La relación de derivación en múltiples pasos se denota por $\Rightarrow^*$.

\begin{definicion}[Lenguaje Generado]
El \textbf{lenguaje generado} por una gramática $G$, denotado $L(G)$, es el conjunto de todas las palabras compuestas exclusivamente por símbolos terminales que pueden derivarse a partir del símbolo inicial $S$:
\[ L(G) = \{ u \in T^* \mid S \Rightarrow^* u \} \].
\end{definicion}

\begin{ejemplo}
Sea la gramática $G = (\{S\}, \{a, b\}, \{S \rightarrow aSb, S \rightarrow \epsilon\}, S)$.
Esta gramática genera el lenguaje $L(G) = \{a^i b^i \mid i \geq 0\}$. Por ejemplo, para generar $aabb$:
$S \Rightarrow aSb \Rightarrow aaSbb \Rightarrow aab\epsilon b = aabb$.
\end{ejemplo}

\subsection{Jerarquía de Chomsky}

Noam Chomsky clasificó las gramáticas en cuatro tipos, estableciendo una jerarquía basada en las restricciones impuestas sobre sus reglas de producción.

\begin{definicion}[Jerarquía de Chomsky]
La jerarquía de Chomsky se define de la siguiente manera:
\begin{itemize}
    \item \textbf{Tipo 0 (Sin restricciones):} Corresponde a cualquier gramática generativa. Estas gramáticas generan los lenguajes \textbf{recursivamente enumerables}, que son aquellos reconocibles por una Máquina de Turing.
    \item \textbf{Tipo 1 (Dependientes del contexto):} Todas las producciones son de la forma $\alpha_1 A \alpha_2 \rightarrow \alpha_1 \beta \alpha_2$, con $A \in V$, $\alpha_1, \alpha_2, \beta \in (V \cup T)^*$ y $\beta \neq \epsilon$. Se permite la regla $S \rightarrow \epsilon$ si $S$ no aparece en la parte derecha de ninguna regla. Generan los lenguajes \textbf{dependientes del contexto}. Un ejemplo es el lenguaje $\{a^n b^n c^n \mid n \geq 1\}$.
    \item \textbf{Tipo 2 (Independientes del contexto):} Todas las producciones son de la forma $A \rightarrow \alpha$, con $A \in V$ y $\alpha \in (V \cup T)^*$. Generan los lenguajes \textbf{independientes del contexto}.
    \item \textbf{Tipo 3 (Regulares):} Todas las producciones son de la forma $A \rightarrow uB$ o $A \rightarrow u$, donde $A, B \in V$ y $u \in T^*$. Generan los \textbf{lenguajes regulares}.
\end{itemize}
\end{definicion}

Esta jerarquía establece una relación de inclusión estricta entre las familias de lenguajes generados, denotadas por $\mathcal{L}_i$ para el tipo $i$:
\[ \mathcal{L}_3 \subset \mathcal{L}_2 \subset \mathcal{L}_1 \subset \mathcal{L}_0 \].
Por ejemplo, todo lenguaje regular es también independiente del contexto, pero no a la inversa.

\begin{thebibliography}{99}

  \bibitem{Referencia1}
  Ismael Sallami Moreno, \textbf{Estudiante del Doble Grado en Ingeniería Informática + ADE}, Universidad de Granada, 2025.
  
  \bibitem{DiapositivasAsignatura}
  Universidad de Granada, \emph{Diapositivas de la asignatura}, Curso 2025/2026.

  % \bibitem{Referencia2}
  % Autor Apellido, \emph{Título del libro o artículo}, Editorial o Revista, Año.
  
  % \bibitem{Referencia3}
  % Nombre Autor, \emph{Título del documento}, Conferencia/URL, Año.
  
  \end{thebibliography}
  

\end{document}
